\subsection{From soft regions to OPE boundary equations}\label{app:boundarymoments}
For a fixed external ordering put \(U=u,\ Y=1-v,\ L=\log U\). A boundary value
\(B_{mNk}=[U^mY^NL^k]I\) is a coefficient of the graph in this chart.
A crossed boundary uses the same definition after relabeling the external
points and transforming the graph normalization. For a trial function
\(\sum_j p_j\sum_\alpha c_{j\alpha} f_\alpha\), one graph coefficient supplies one equation,
\[
 \sum_{j,\alpha}c_{j\alpha}[U^mY^NL^k]\bigl(p_j f_\alpha\bigr)=B_{mNk}.
\]
Expand the complete product. A pole in \(p_j\) can bring a higher power of a
word into the requested order. In the radial chart \(z=th,\ \bar z=t/h\) one
has \(U=t^2,\ Y=t(h+h^{-1})-t^2\). Thus a coefficient of \(t^nL^k\) at
fixed angle is a specified linear combination of the \(B_{mNk}\). This is how
the graph moments and the svMPL functions enter the same linear system.

\paragraph{Calculating the graph side.}
Fix \(x_1=0,\ x_3=e,\ e^2=1,\ x_4=\infty\), and take \(x_2=q\) near zero,
with \(q^2=U\) and \(2e\cdot q=U+Y\). In \(D=4-2\epsilon\) dimensions, assign
each integration point to the soft region (distance of order \(|q|\) from
zero) or the hard region (distance of order one). Expand each mixed propagator
as
\[
 \frac1{(x-\xi)^2}=\frac1{x^2}
 \sum_{\ell\ge0}\left(\frac{2x\cdot\xi-\xi^2}{x^2}\right)^\ell ,
 \qquad x\ {\rm hard},\quad \xi\ {\rm soft}.
\]
Keep soft and hard subgraph propagators exact. Integrate the resulting tensors in
each subgraph, retain the external normalization, and sum regions before taking
\(\epsilon\to0\). With \(s\) soft loops the scale factor is \(U^{-s\epsilon}\)
after extracting the integer power of \(U\). Its expansion generates the OPE
logarithms. For example, if \(a_s\) is the total quartic-pole coefficient at
fixed \(m,N\), finiteness gives \(\sum_s s^j a_s=0\) for \(j=0,1,2,3\).
Consequently \((a_0,\ldots,a_4)=A(1,-4,6,-4,1)\) and
\(B_{mN4}=\sum_s s^4a_s/4!=a_2/6\). Lower logarithms require the appropriate
subleading Laurent coefficients.

\needspace{9\baselineskip}
\paragraph{Two held-out moments of \(I_{212}\).}
Here \(C_N=[U^0Y^NL^2]I_{212}\) in its primary OPE ordering.
The fitted primary moments end at \(N=19\). In the two-soft calculation
relabel the external vertices of basis graph 212 by \((2,4,1,3)\).
Vertices \(6,8\) are soft and \(5,7\) are hard. Project onto the leading
traceless part by taking \(q^2=0\) in the Taylor tensors, after the scale
logarithm has been extracted. Invert the hard coordinates, denoting them
by \(x,y\), and put \(X=2q\cdot x,\ Z=2q\cdot y\).
Write \(a=\alpha,\ b=\alpha+(1-\alpha)\tau\). The two-soft residue acts on
angular monomials as
\[
 S_{rs}=\frac12\int_0^1\!\mathrm d\alpha\int_0^1\!\mathrm d\tau\,a^r b^s
       =\frac1{2(s+1)}\sum_{\ell=0}^{s}\frac1{r+\ell+1}.
\]
The mixed denominators \(s_{25},s_{58},s_{67}\) give the three geometric
series in the following expression. The graph numerator \(s_{27}s_{56}\)
supplies the polynomial factors in
\[
 K(X,Z;a,b)=\frac{(1-Z)(1-aX)}{(1-X)(1-bX)(1-aZ)} .
\]
For \(i+j=N\), replace \(a^r b^s\) in \([X^iZ^j]K\) by \(S_{rs}\), and denote
the result by \(\omega_{ij}\). The remaining integral \(H_{ij}\) is the finite
hard two-loop propagator moment with tensor \(X^iZ^j\), normalized to
\(2q\cdot e=1\). With \(A=x^2,B=(x-e)^2,C=y^2,D'=(y-e)^2,E=(x-y)^2\),
its denominators are \(ABCD'E\). Integrating the total derivative with
vector \(x-y\) in \(x\) gives
\[
 (D-4+i)H_{ij}-iH_{i-1,j+1}
 =\int\frac{\mathrm d^Dx\,\mathrm d^Dy}{\pi^D}
 \frac{X^i Z^j}{ABCD'E}\left(\frac{E-C}{A}+\frac{E-D'}{B}\right).
\]
The right side reduces to products of one-loop bubbles. At \(i=0\) keep the
order-\(\epsilon\) term since \(D-4=-2\epsilon\). Then step up in \(i\) at
fixed \(i+j\). The sum \(C_N=\sum_{i+j=N}\omega_{ij}H_{ij}\) gives
\begin{align*}
 C_{22}&=\frac{7436781516927153228656542409054377847}{58972189577941010378035403517542400000}
       +\frac{1332950097}{62941935056}\,\zeta(3),\\
 C_{23}&=\frac{1319341114152286490801205455383238191}{10701947447036553868225706498457600000}
       +\frac{1347822955}{68534129664}\,\zeta(3).
\end{align*}
The bootstrap comparison expands the frozen HPL expression at \(Y=0\) with
kernels \(\mathrm dY/Y,\mathrm dY/(Y-1)\). Its rational and \(\zeta(3)\)
parts agree separately at both orders. The same graph sum agrees through all
orders \(0\le N\le23\). The \href{https://github.com/windfolgen/svbwalkthrough/tree/main/runs/fourloopI212}{boundary calculation}
contains the two independent evaluations.
